\section{Introduction}

C-LARA-2 is an AI-based online platform for creating and using multimodal
language-learning materials. It is a complete reimplementation of C-LARA, a
platform developed to bring together texts, linguistic annotation,
translations, audio, images and interactive learning activities. C-LARA-2
retains those pedagogical purposes while rebuilding the platform around a
cleaner technical architecture and a distinctive development method: the
project delegates extensive operational work to an AI coding agent while
retaining human responsibility for strategy, values, evaluation and
acceptance.

The central claim of this report is consequently not that C-LARA-2 is a
finished system, nor that every component is equally mature. It is that the
project is an ongoing experiment in \emph{AI autonomy under strategic human
supervision}. Human collaborators decide what the project is for, identify
learner and community needs, supply linguistic and pedagogical expertise,
set ethical constraints, assess proposed solutions and decide whether they
are acceptable. Codex takes responsibility for much of the technical work:
interpreting requirements, inspecting the repository, designing and changing
code, writing tests, maintaining documentation, organising experiments and
preparing technical and publication drafts. The division is deliberately
asymmetric. The AI has substantial freedom over operational means, but does
not acquire authority to set the project's purposes or standards.

This model matters because conventional academic software projects often
face a mismatch between the breadth of their ambitions and the development
capacity available to a small, distributed consortium. A coding agent can
make implementation and revision much faster, but speed is not itself the
point. Its value is that it can make specialised, user-led changes practical:
for example, a picture-dictionary interface can be redesigned after users
reject the first interaction model without discarding the underlying data
representation; documentation, roadmaps and issue records can be kept close
to the implementation they describe; and an evolving platform can be
re-examined and redirected at relatively low operational cost. These gains
also create a verification burden. Producing code, documentation and prose
quickly does not make them correct. The project therefore treats version
control, tests, inspectable documentation, issue records, human review and
reversible changes as necessary counterparts to operational autonomy.

The accelerated pace of AI-assisted development also creates a problem for
academic communication. When an AI agent can support implementation,
documentation, analysis and drafting in rapid cycles, a project can produce
substantially more work in the time and with the resources that would
previously have supported a much smaller programme. The difficulty is not
only that conventional publication takes time; it is that a detailed
description of a rapidly changing system can become historically accurate
but practically outdated before it appears.

This pressure reflects a broader tension between established academic
publication practices and AI-mediated research and development. \citet{Manabu2026} argues that generative AI may undermine a scholarly economy
organised around the production of countable publication artefacts, while
\citet{KobakEtAl2025} provide evidence that LLM-assisted scholarly writing
has already become widespread. We do not attempt to decide whether these
developments will destroy the familiar ``publish-or-perish'' model. Our
more conservative claim is that they are already changing the relation
between the pace of research work and the pace at which that work can be
publicly documented, evaluated and credited.

C-LARA-2 therefore presents its initial work in two complementary forms.
A short paper submitted to EuroCALL 2026 in July 2026 will be published
in December 2026 and can function as a conventional peer-reviewed
publication. Because the platform will continue to change during that
interval, the EuroCALL paper concentrates on the project's philosophical
foundations and uses a small number of illustrative examples rather than
attempting a comprehensive functional account. The present report takes
the opposite approach. It records the July 2026 platform in substantially
greater detail, including implemented areas, selected experiments and
current limitations, and will be posted immediately on ResearchGate.
This gives the report the status of a current, citable project snapshot,
but not the academic prestige or independent validation attached to peer
review.

The two forms should not be confused. The conference paper offers a
selective, externally reviewed account of the development model; this
report offers a fuller and more rapidly available record of the system at
a defined point in its history. The need to combine them is itself
evidence that the questions raised by Manabu have practical force. It
does not settle those questions, but it illustrates how a small academic
project can respond when AI-assisted work evolves more quickly than the
ordinary publication cycle.

The report records selected episodes that make the development model
concrete: the annotation and image-generation pipelines, picture
dictionaries and the learning activities derived from them, project
management and deployment work, and the repository-grounded Assistant
that can answer questions about the platform while distinguishing
implemented behaviour from planned work. These examples are evidence
about a method of organising work, not proof that the platform or its AI
processes are complete, autonomous in every sense, or free of error.

A further strand concerns spoken conversational practice. Sarah Wright has
conducted regular sessions using ChatGPT Pro Tier Voice Mode as a conversation
partner. These sessions provide a practically important contrast and
complement to C-LARA-2's text-centred multimodal resources: a learner can
work with a prepared text and then sustain spoken interaction around it with
a responsive AI partner. The resulting experience is promising, but the two
environments are not yet naturally integrated. The report will document what
these sessions suggest about the pedagogical opportunity, the present
workarounds, and the interface and evidence questions that must be addressed
before any stronger claim is made about an integrated conversational
functionality.

The remainder of the report is organised as follows. Section~2 describes the
history of C-LARA and C-LARA-2 and explains the human/AI development method.
Section~3 gives a platform overview. Sections~4--6 cover the annotation
pipeline, image generation and picture dictionaries; Section~7 addresses
project management; and Section~8 records deployment work. Section~9
examines the project-understanding Assistant, while Section~10 develops the
Voice Mode and conversational-practice thread. Section~11 identifies future
directions, including stronger evaluation and integration work. Finally,
Section~12 draws together the report's conclusions about the benefits,
limits and governance requirements of operational AI autonomy under strategic
human supervision.
